\chapter{Reference: HTTP and tasks}
\label{ch:refman-task}

\begin{aside}
Anything that blocks belongs in a command. Files, network,
process spawns — return a \code{task} or \code{http\_get}
and \maya{} will run it off the loop, then post the result
back as a message.
\end{aside}

\section{Forms}

\begin{itemize}[leftmargin=*]
  \item \code{Cmd<Msg>::task(callable, then\_msg\_fn)}: run
    \code{callable()} on a worker, pass the result to
    \code{then\_msg\_fn} to make a \code{Msg}.
  \item \code{Cmd<Msg>::http\_get(url, then\_msg\_fn)}: GET
    request, dispatch a message with the body or error.
  \item \code{Cmd<Msg>::read\_file(path, then\_msg\_fn)}:
    async file read.
\end{itemize}

\section{Working example}

\begin{lstlisting}[language=C++]
#include <maya/maya.hpp>
#include <thread>

using namespace maya;
using namespace maya::dsl;

struct TaskDemo {
    struct Model {
        std::string status = "idle";
        long answer = 0;
    };
    struct Start {};
    struct Done { long value; };
    struct Quit {};
    using Msg = std::variant<Start, Done, Quit>;

    static Model init() { return {}; }

    static auto update(Model m, Msg msg)
        -> std::pair<Model, Cmd<Msg>>
    {
        return std::visit(overload{
            [&](Start) {
                m.status = "computing...";
                auto compute = []() -> long {
                    std::this_thread::sleep_for(
                        std::chrono::seconds(2));
                    long sum = 0;
                    for (long i = 1; i <= 1000000; ++i) sum += i;
                    return sum;
                };
                return std::pair{m,
                    Cmd<Msg>::task(compute,
                        [](long v) { return Done{v}; })};
            },
            [&](Done d) {
                m.status = "done";
                m.answer = d.value;
                return std::pair{m, Cmd<Msg>{}};
            },
            [](Quit) {
                return std::pair{Model{}, Cmd<Msg>::quit()};
            },
        }, msg);
    }

    static Element view(const Model& m) {
        return v(
            text(m.status) | Bold,
            text(std::format("answer: {}", m.answer)),
            blank_,
            t<"s start, q quit"> | Dim
        ) | pad<1> | border_<Round>;
    }

    static auto subscribe(const Model&) -> Sub<Msg> {
        return key_map<Msg>({
            {'q', Quit{}},
            {'s', Start{}},
        });
    }
};

int main() { run<TaskDemo>({.title = "task"}); }
\end{lstlisting}

\section{Notes}

\begin{itemize}[leftmargin=*]
  \item The worker pool is shared; long tasks block other
    tasks, not the UI.
  \item Exceptions in the callable become a \code{Done}
    with a default-constructed value plus an error flag if
    your \code{Msg} carries one.
  \item Never touch the Model from inside the task. Only
    the returned message may mutate state.
\end{itemize}

\section{Recap}

\begin{recap}
\begin{itemize}[leftmargin=*]
  \item \code{task}, \code{http\_get}, \code{read\_file}.
  \item Blocking work goes off-loop; result returns as a
    message.
  \item No shared state with the worker.
\end{itemize}
\end{recap}

\section*{Workshop}

\begin{enumerate}[leftmargin=*]
  \item Spawn five tasks at once; show progress as each
    finishes.
  \item Fetch a URL and render the first 200 bytes.
  \item Read a file and count its lines off-loop.
\end{enumerate}
